
\documentclass{article}
\usepackage{colortbl}
\usepackage{makecell}
\usepackage{multirow}
\usepackage{supertabular}

\begin{document}

\newcounter{utterance}

\centering \large Interaction Transcript for game `cladder', experiment `full\_v1.5\_default', episode 7517 with Qwen3.5{-}9B.
\vspace{24pt}

{ \footnotesize  \setcounter{utterance}{1}
\setlength{\tabcolsep}{0pt}
\begin{supertabular}{c@{$\;$}|p{.15\linewidth}@{}p{.15\linewidth}p{.15\linewidth}p{.15\linewidth}p{.15\linewidth}p{.15\linewidth}}
    \# & \multicolumn{2}{c}{Player} && \multicolumn{2}{c}{Game Master} \\
    \hline

    \theutterance \stepcounter{utterance}  
    & & & \multicolumn{4}{p{0.6\linewidth}}{
        \cellcolor[rgb]{0.9,0.9,0.9}{
            \makecell[{{p{\linewidth}}}]{
                \texttt{\tiny{[P1$\langle$GM]}}
                \texttt{You are an expert in causal inference. The following question is not a typical commonsense query, but rather a meticulously designed question created by a professor specializing in causal inference, intended to assess the students' mastery of the course content.} \\
\texttt{Imagine a self{-}contained, hypothetical world with only the following conditions, and without any unmentioned factors or causal relationships: Wibl has a direct effect on vubr and zuph. Zuph has a direct effect on uvzi. Vubr has a direct effect on uvzi. Method 1: We look at how zuph correlates with uvzi case by case according to vubr. Method 2: We look directly at how zuph correlates with uvzi in general. To understand how zuph affects uvzi, is it more correct to use the Method 1 than Method 2?} \\
\texttt{Start your answer with 'Yes' or 'No', followed by additional reasoning or evidence to support your explanation.} \\
            }
        }
    }
     \\ \\

    \theutterance \stepcounter{utterance}  
    & \multicolumn{4}{p{0.6\linewidth}}{
        \cellcolor[rgb]{0.561,0.737,0.561}{
            \makecell[{{p{\linewidth}}}]{
                \texttt{\tiny{[P1$\rangle$GM]}}
                \texttt{No} \\
\\ 
\texttt{**Reasoning:**} \\
\\ 
\texttt{In this hypothetical world, the causal structure forms a **confounding scenario** (specifically, a "fork" or common cause structure relative to the variables of interest, though here it is a collider{-}like dependency in terms of correlation patterns).} \\
\\ 
\texttt{1.  **Causal Graph Analysis**:} \\
\texttt{    *   \$Wibl \textbackslash{}rightarrow Zuph\$} \\
\texttt{    *   \$Wibl \textbackslash{}rightarrow Vubr\$} \\
\texttt{    *   \$Zuph \textbackslash{}rightarrow Uvzi\$} \\
\texttt{    *   \$Vubr \textbackslash{}rightarrow Uvzi\$} \\
\\ 
\texttt{    Here, \$Wibl\$ is a common cause of both \$Zuph\$ and \$Vubr\$. Consequently, \$Zuph\$ and \$Vubr\$ are correlated due to their shared cause (\$Wibl\$). Since both \$Zuph\$ and \$Vubr\$ independently affect the outcome \$Uvzi\$, the total correlation between \$Zuph\$ and \$Uvzi\$ (Method 2) is a mixture of the direct causal effect of \$Zuph\$ and the spurious correlation induced by \$Wibl\$ (which correlates \$Zuph\$ with \$Vubr\$, which in turn affects \$Uvzi\$).} \\
\\ 
\texttt{2.  **Evaluating Method 2 (Unconditional Correlation)**:} \\
\texttt{    Method 2 calculates the marginal correlation between \$Zuph\$ and \$Uvzi\$. Because \$Zuph\$ and \$Vubr\$ are correlated (via \$Wibl\$), and \$Vubr\$ causes \$Uvzi\$, the observed correlation between \$Zuph\$ and \$Uvzi\$ will be biased. It will overestimate (or underestimate, depending on the signs of the effects) the true causal effect of \$Zuph\$ on \$Uvzi\$ because it conflates the direct path (\$Zuph \textbackslash{}rightarrow Uvzi\$) with the indirect path through the confounder (\$Zuph \textbackslash{}leftarrow Wibl \textbackslash{}rightarrow Vubr \textbackslash{}rightarrow Uvzi\$).} \\
\\ 
\texttt{3.  **Evaluating Method 1 (Stratified Analysis)**:} \\
\texttt{    Method 1 examines the correlation between \$Zuph\$ and \$Uvzi\$ "case by case according to \$Vubr\$." This is equivalent to **stratifying** on \$Vubr\$ (or conditioning on it).} \\
\texttt{    *   By holding \$Vubr\$ constant (fixing the value of \$Vubr\$), we block the backdoor path \$Zuph \textbackslash{}leftarrow Wibl \textbackslash{}rightarrow Vubr \textbackslash{}rightarrow Uvzi\$.} \\
\texttt{    *   Once \$Vubr\$ is controlled for, the remaining variation in \$Uvzi\$ is no longer influenced by the shared cause \$Wibl\$ acting through \$Vubr\$.} \\
\texttt{    *   Therefore, the correlation observed within each stratum of \$Vubr\$ isolates the direct causal} \\
            }
        }
    }
    & & \\ \\

    \theutterance \stepcounter{utterance}  
    & & & \multicolumn{2}{p{0.3\linewidth}}{
        \cellcolor[rgb]{0.9,0.9,0.9}{
            \makecell[{{p{\linewidth}}}]{
                \texttt{\tiny{[GM$|$GM]}}
                \texttt{No} \\
            }
        }
    }
    & & \\ \\

    \theutterance \stepcounter{utterance}  
    & & & \multicolumn{2}{p{0.3\linewidth}}{
        \cellcolor[rgb]{0.9,0.9,0.9}{
            \makecell[{{p{\linewidth}}}]{
                \texttt{\tiny{[GM$|$GM]}}
                \texttt{yes} \\
            }
        }
    }
    & & \\ \\

    \theutterance \stepcounter{utterance}  
    & & & \multicolumn{2}{p{0.3\linewidth}}{
        \cellcolor[rgb]{0.9,0.9,0.9}{
            \makecell[{{p{\linewidth}}}]{
                \texttt{\tiny{[GM$|$GM]}}
                \texttt{game\_result = LOSE} \\
            }
        }
    }
    & & \\ \\

\end{supertabular}
}

\end{document}
